\section*{Capítulo \ref{ch:b2:current-potential-curves} --- Curvas corrente--potencial}

\begin{solution}{exo:b2:current-potential-curves:1}
O zinco é oxidado: \omterm{def:b2:current-potential-curves:ie-curve}{corrente anódica}, positiva. Os íons cobre são reduzidos no
eletrodo de cobre: \omterm{def:b2:current-potential-curves:ie-curve}{corrente catódica}, negativa. As duas correntes são iguais
em valor absoluto, pois a mesma carga circula pelo circuito.
\end{solution}

\begin{solution}{exo:b2:current-potential-curves:2}
$|i_{\lim}| = 96\,485 \times 0.50 \times 10^{-4} \times 1.6 \times 10^{-9} \times
2.0/(25 \times 10^{-6}) = \qty{6.2e-4}{A}$, isto é, \qty{0.62}{mA}.
\end{solution}

\begin{solution}{exo:b2:current-potential-curves:3}
Um \omterm{def:b2:current-potential-curves:fast-slow}{sistema rápido} corta o eixo dos potenciais abruptamente em seu potencial de
Nernst; um \omterm{def:b2:current-potential-curves:fast-slow}{sistema lento} apresenta um trecho plano de corrente nula em torno
dele, e seus ramos só começam após uma \omterm{def:b2:current-potential-curves:overpotential}{sobretensão}. Sobre a platina:
\ce{Fe^3+}/\ce{Fe^2+} (rápido); \ce{O2}/\ce{H2O} (lento).
\end{solution}

\begin{solution}{exo:b2:current-potential-curves:4}
No potencial de meia-onda, igual a $E^\circ = \qty{0.77}{V}$ quando as duas
formas difundem do mesmo modo.
\end{solution}

\begin{solution}{exo:b2:current-potential-curves:5}
Perto de \qty{0.77}{V}, uma onda catódica do \ce{Fe^3+} (rápido), que atinge um
patamar proporcional à sua concentração. Perto de \qty{0.34}{V}, uma segunda
onda, com deposição de cobre; seu patamar se soma ao primeiro e é cerca de
duas vezes mais alto para uma concentração igual e coeficientes de difusão
semelhantes, já que $n = 2$. Depois, perto de \qty{0}{V} a pH 0 sobre a platina
(mais abaixo quando o eletrodo fica recoberto de cobre), a barreira abrupta
da redução de \ce{H+}.
\end{solution}

\begin{solution}{exo:b2:current-potential-curves:6}
A \qty{0.50}{V}: só o \ce{Fe^3+} é reduzido a \ce{Fe^2+}, em sua \omterm{def:b2:current-potential-curves:limiting-current}{corrente limite}.
A \qty{0.10}{V}: o \ce{Fe^3+} é reduzido e o cobre se deposita ao mesmo tempo,
cada um em sua \omterm{def:b2:current-potential-curves:limiting-current}{corrente limite}; o total é a soma dos dois patamares.
\end{solution}

\begin{solution}{exo:b2:current-potential-curves:7}
pH 0: \qty{0}{V} e \qty{1.23}{V}; pH 14: \qty{-0.83}{V} e \qty{0.40}{V}. Sobre a
platina, a janela mantém sua largura e desce \qty{59}{mV} por unidade de pH,
cada barreira acompanhando seu par, a anódica ainda deslocada para cima pela
\omterm{def:b2:current-potential-curves:overpotential}{sobretensão} do oxigênio.
\end{solution}

\begin{solution}{exo:b2:current-potential-curves:8}
A \qty{1.0}{V}, o \ce{Fe^2+} é oxidado em seu patamar (\omterm{def:b2:current-potential-curves:ie-curve}{corrente anódica}
proporcional a $[\ce{Fe^2+}]$) e todo \ce{Ce^4+} presente é reduzido em seu
patamar (\omterm{def:b2:current-potential-curves:ie-curve}{corrente catódica} proporcional a $[\ce{Ce^4+}]$). Antes da
equivalência, a \omterm{def:b2:current-potential-curves:ie-curve}{corrente anódica} cai linearmente com o volume adicionado;
depois dela, a \omterm{def:b2:current-potential-curves:ie-curve}{corrente catódica} cresce linearmente. As duas retas cortam o
zero na equivalência.
\end{solution}

\begin{solution}{exo:b2:current-potential-curves:9}
Os patamares dobram ($|i_{\lim}| \propto 1/\delta$). O potencial de meia-onda não
muda, já que a camada é a mesma para as duas formas; nem o potencial de
corrente nula de um \omterm{def:b2:current-potential-curves:fast-slow}{sistema rápido}, que é o potencial de Nernst.
\end{solution}

\begin{solution}{exo:b2:current-potential-curves:10}
Como no capítulo, com $i_{\lim,c} = 0$: $E = E_{1/2} + (RT/nF)\ln[i/(i_{\lim} -
i)]$. Em logaritmos decimais, $E = E_{1/2} + (RT\ln 10/nF)\log[i/(i_{\lim} - i)]$,
uma reta de inclinação $0.0592/n$ volt a \qty{298}{K}, que corta o eixo de
$E$ em $E_{1/2}$.
\end{solution}

\begin{solution}{exo:b2:current-potential-curves:11}
Sozinho, o zinco se estabiliza no potencial em que sua \omterm{def:b2:current-potential-curves:ie-curve}{corrente anódica}
(oxidação rápida do zinco) é igual à \omterm{def:b2:current-potential-curves:ie-curve}{corrente catódica} de redução de \ce{H+}
sobre o zinco, que é pequena porque essa redução é lenta sobre o zinco: a
dissolução é lenta. Em contato com a platina, os elétrons liberados pelo zinco
podem reduzir \ce{H+} sobre a platina, onde a redução é rápida: a \omterm{def:b2:current-potential-curves:ie-curve}{corrente
catódica} a um dado potencial é muito maior, o equilíbrio de correntes é
atingido a uma corrente mais alta, e o zinco se dissolve mais depressa
enquanto o hidrogênio se forma sobre a platina.
\end{solution}

\begin{solution}{exo:b2:current-potential-curves:12}
Sobre a platina, a oxidação da água começa perto de \qty{1.6}{V} a pH 0: antes
de \qty{2.0}{V}, a corrente seria gasta produzindo dioxigênio. Um eletrodo sobre
o qual a oxidação da água exige uma \omterm{def:b2:current-potential-curves:overpotential}{sobretensão} maior empurra a barreira para
além de \qty{2.0}{V}, de modo que a espécie pode ser oxidada dentro da janela.
\end{solution}

\begin{solution}{pb:b2:current-potential-curves:1}
\textbf{1.} \ce{O2 + 2H2O + 4e- -> 4OH-}.
\textbf{2.} \ce{Ag + Cl- -> AgCl + e-} (quatro vezes); global \ce{O2 + 2H2O + 4Ag +
4Cl- -> 4AgCl + 4OH-}.
\textbf{3.} $E = 1.229 - 0.0592 \times 7 + (0.0592/4)\log 0.21 = \qty{0.80}{V}$.
\textbf{4.} $E^\circ(\ce{AgCl}/\ce{Ag}) = \qty{0.22}{V}$ com cloreto de atividade 1:
$-0.60 + 0.22 = \qty{-0.38}{V}$ na escala do hidrogênio.
\textbf{5.} A redução do \ce{O2} é um \omterm{def:b2:current-potential-curves:fast-slow}{sistema lento}: abaixo de seu potencial de
Nernst, é necessária uma \omterm{def:b2:current-potential-curves:overpotential}{sobretensão} de vários décimos de volt antes que a
corrente cresça, e mais ainda para atingir o patamar.
\textbf{6.} No patamar, a corrente depende apenas do suprimento de \ce{O2},
e não de pequenas derivas do potencial: ela mede a concentração.
\textbf{7.} A redução da água a di-hidrogênio (a barreira catódica), que
acrescentaria uma corrente sem relação com o oxigênio.
\textbf{8.} Linear através da membrana, de $c$ do lado da amostra até zero no
cátodo.
\textbf{9.} $j = D_m c/L$ por unidade de área.
\textbf{10.} $i = -4FAD_m c/L$ (catódica), $|i| = 4FAD_mc/L$.
\textbf{11.} A membrana é a etapa mais lenta: a difusão através dela é muito
mais lenta que na água do lado de fora; por isso a membrana faz o papel da
\omterm{def:b2:current-potential-curves:limiting-current}{camada de difusão}.
\textbf{12.} O gradiente de concentração fica dentro da membrana, cuja
espessura não depende da agitação (desde que a água do lado de fora se
renove); um depósito sobre a membrana aumenta a espessura e diminui a corrente.
\textbf{13.} $0.2095 \times (101.325 - 3.17) = \qty{20.56}{kPa}$.
\textbf{14.} $c = 1.3 \times 10^{-5} \times 20\,560 = \qty{0.27}{mol/m^3}$ ($0.267$).
\textbf{15.} $0.267 \times 32.00 = \qty{8.6}{mg/L}$ (\unit{g/m^3}).
\textbf{16.} $|i| = 4 \times 96\,485 \times 1.0 \times 10^{-5} \times 1.0 \times 10^{-10}
\times 0.267/(25 \times 10^{-6}) = \qty{4.1}{\micro A}$.
\textbf{17.} $4.05/8.55 = \qty{0.474}{\micro A}$ per \unit{mg/L}.
\textbf{18.} A espessura da membrana e seu coeficiente de difusão não são
conhecidos com precisão e mudam com o envelhecimento e com a temperatura; uma
medida em uma solução conhecida inclui todos eles.
\textbf{19.} Zero (na prática, uma pequena corrente residual, que é subtraída).
\textbf{20.} $2.80/0.474 = \qty{5.9}{mg/L}$.
\textbf{21.} $2.80/4.05 = 69~\%$ da saturação.
\textbf{22.} Oxigênio consumido pela decomposição de matéria orgânica (esgoto,
plantas mortas) mais depressa do que é reposto pelo ar e pela fotossíntese.
\textbf{23.} $2.80 \times 10^{-6}/(4 \times 96\,485) \times 3600 = \qty{2.6e-8}{mol}$
por hora: desprezível diante do oxigênio de um único litro de água
(\qty{1.8e-4}{mol}).
\textbf{24.} Tanto a solubilidade do \ce{O2} (\omterm{prop:b2:chemical-potential:henry}{constante de Henry}) quanto a
difusão através da membrana variam com a temperatura: a sensibilidade
determinada a \qty{25}{\degreeCelsius} deixa de valer.
\textbf{25.} A amostra contém $\boldsymbol{\approx \qty{5.9}{mg/L}}$ de oxigênio
dissolvido.
\end{solution}
